\documentclass[../main.tex]{subfiles}

\begin{document}

\chapter{测度与拓扑}\label{chapter:ch21-测度与拓扑}

在欧氏空间 $\mathbf{R}^{n}$ 特别是实直线上的 Lebesgue 测度 $\mu_{n}$ 与 Lebesgue 积分的研究中, 我们考察了 Lebesgue 测度与拓扑以及可测函数与连续函数的联系. Borel $\sigma$ 代数 $\mathcal{B}(\mathbf{R}^{n})$ 是包含于 Lebesgue 可测集的代数. 因此, 若我们定义 $C_{c}(\mathbf{R}^{n})$ 为 $\mathbf{R}^{n}$ 上在闭有界集外消失的连续实值函数的线性空间, 算子
\begin{align*}
\text{对所有 }f\in C_{c}(\mathbf{R}^{n}),\quad f\mapsto\int_{\mathbf{R}^{n}}f\,d\mu_{n}
\end{align*}
是恰当定义的、正的与线性的. 此外, 对 $\mathbf{R}^{n}$ 的闭有界子集 $K$, 算子
\begin{align*}
\text{对所有 }f\in C_{c}(K),\quad f\mapsto\int_{K}f\,d\mu_{n}
\end{align*}
是恰当定义的、正的, 且若 $C(K)$ 有最大值范数, 它是有界线性算子.

本章我们考虑一般的局部紧拓扑空间 $(X,\mathcal{T})$、由包含拓扑 $\mathcal{T}$ 的最小 $\sigma$ 代数组成的 Borel $\sigma$ 代数 $\mathcal{B}(X)$、关于 Borel 测度 $\mu:\mathcal{B}(X)\to[0,\infty)$ 的积分. 本章有两个中心. 第一个是 Riesz-Markov 定理, 它告诉我们, 若 $C_{c}(X)$ 表示在紧集外消失的 $X$ 上连续实值函数的线性空间, 则 $C_{c}(X)$ 上的每个正线性泛函由关于 $\mathcal{B}(X)$ 上的 Borel 测度的积分给出. Riesz-Markov 定理使我们能够证明 Riesz 表示定理, 它告诉我们, 对紧 Hausdorff 拓扑空间 $X$, 赋予最大值范数的线性空间 $C(X)$ 上的每个有界线性泛函由关于带号 Borel 测度的积分给出.

此外, 在这些表示定理中可以选取表示测度属于一类 Borel 测度, 我们把它命名为 Radon 测度, 在该类中表示测度是唯一的. Riesz 表示定理提供了将 Alaoglu 定理、Helly 定理和 Krein-Milman 定理应用于 Borel 测度族的机会.

这两个表示定理的证明需要探讨集合上的拓扑以及与该拓扑相关的 Borel 集上的测度之间的关系. 我们构造表示泛函的 Borel 测度所用的技巧与构造欧氏空间上的 Lebesgue 测度所用的技巧相同: 研究定义在 $X$ 的特定子集族 $\mathcal{S}$ 上的预测度的 Carath\'eodory 延拓, 现在取 $\mathcal{S}=\mathcal{T}$, 即 $X$ 上的拓扑. 本章我们从局部紧拓扑空间的预备知识开始. 21.2 节讨论将这类空间的所有性质写成一个定理并且对 $X$ 是局部紧度量空间的情形给出该定理非常简单的证明.

\subfile{Chapter-21/section-01.tex}
\subfile{Chapter-21/section-02.tex}
\subfile{Chapter-21/section-03.tex}
\subfile{Chapter-21/section-04.tex}
\subfile{Chapter-21/section-05.tex}
\subfile{Chapter-21/section-06.tex}

\end{document}